\begin{textsample}{POS Dim 2 – 2014\_c5 – Score -1.00 – 2014\_c5/enade\_2014\_pedag\_076bc59f40.txt}  \label{ex:f2_pos_310}
CAPÍTULO 5 DISTRIBUIÇÃO DOS CONCEITOS 5.1 PANORAMA NACIONAL DA DISTRIBUIÇÃO DOS CONCEITOS O Conceito Enade é um indicador de qualidade que avalia o desempenho dos estudantes a partir dos resultados do ENADE.
Os valores possíveis vão de 1 ( pior situação ) a 5 ( melhor situação ).
As Unidades de Observação com apenas um ou sem nenhum concluinte participante não obtêm o Conceito Enade, ficando Sem Conceito ( SC ).
A Tabela 5.1 apresenta a quantidade e distribuição de cursos de Pedagogia ( Licenciatura ) participantes do ENADE/2014, por faixa de conceito e Grande Região.
Mais uma vez, ressaltando que a diferença entre os cursos tabulados neste capítulo e no capítulo 2 são os cursos sem conceito, em princípio, aqueles sem alunos concluintes que participassem da prova ou que tivessem tão somente um aluno concluinte21.
Observando -se os dados da Tabela 5.1, nota -se que, dos 1.126 cursos participantes, 481 ( 42,7 \% ) classificaram -se com conceito 3, o valor modal.
Este foi também o conceito modal em mais três regiões : Nordeste ( 42,4 \% ), Sudeste ( 46,9 \% ) e Sul ( 41,3 \% ).
Nas regiões Norte e Centro-Oeste, a moda foi o conceito 2 ( 48,9 \% e 39,1 \%, respectivamente ).
O conceito 2 foi o segundo mais frequente em nível nacional ( 24,5 \%, correspondendo a 276 cursos ) e o conceito 4, o terceiro ( 23,4 \%, correspondendo a 264 cursos ).
Houve, ainda, 35 cursos ( 3,1 \% ) que receberam conceito 5 e 47 ( 4,2 \% ) cursos que receberam conceito 1.
Dos 1.126 cursos de Pedagogia ( Licenciatura ), 23 cursos ( 2,0 \% ) ficaram sem conceito ( SC ). 21 Estes cursos com somente um concluinte não foram considerados no capítulo 2.
A região Norte participou com 94 cursos ou 8,3 \% do total nacional.
Nesta região, o conceito 2 foi atribuído a 46 cursos dos 94 participantes, o que equivale a 48,9 \% do total regional, além de ser o conceito modal na região.
Nenhum curso ficou com conceito 5 e três cursos ( 3,2 \% ) ficaram sem conceito ( SC ).
Enquanto, 14 cursos foram avaliados com conceitos 1 e 4 ( sete cursos cada ), correspondendo a 7,4 \% cada e 31 cursos ( 33,0 \% ) com conceito 3.
A região Nordeste participou com 238 cursos ou 21,1 \% do total nacional.
Nesta região, 101 cursos, 42,4 \% em termos regionais, obtiveram conceito 3, o conceito modal para a região.
Os conceitos 2 e 4 foram atribuídos a, respectivamente, 72 e 30 cursos ( 30,3 \% e 12,6 \% ), e 27 cursos ( 11,3 \% ) receberam conceito 1.
Oito cursos obtiveram o conceito 5 ou ficaram sem conceito ( SC ) ( quatro cursos, cada ), correspondendo a 1,7 \% cada, nessa região.
Dos 495 cursos participantes da região Sudeste ( 44,0 \% do total nacional ), o conceito 3 também foi o conceito modal, tendo sido obtido por 232 cursos ( 46,9 \% ).
O conceito 1 foi atribuído a seis cursos ( 1,2 \% ) e os conceitos 2 e 4, a 91 ( 18,4 \% ) e 135 ( 27,3 \% ) cursos, respectivamente.
O conceito 5 foi atribuído a 21 cursos ( 4,2 \% ).
E dez cursos ( 2,0 \% ) ficaram sem conceito ( SC ).
Os 189 cursos da região Sul ficaram distribuídos em quase todas as faixas de conceitos, exceto no conceito 1.
Há predominância do conceito 3 com 41,3 \% ( conceito modal ), ou seja, 78 dos 189 cursos participantes na região Sul.
O conceito 2 foi atribuído a 24 cursos ( 12,7 \% ) e o conceito 4, a 72 cursos ( 38,1 \% ).
O conceito 5 foi recebido por nove cursos ( 4,8 \% ).
Seis cursos na região Sul ficaram sem conceito ( SC ).
Já dos 110 cursos participantes na região Centro-Oeste, 43 receberam conceito 2, o conceito modal ( 39,1 \% ).
Nesta região, nenhum curso ficou sem conceito ( SC ).
Já 39 cursos foram alocados no conceito 3, correspondendo a 35,5 \% dos cursos nessa região.
Sete cursos ( 6,4 \% ) foram avaliados com conceito 1, 20 cursos ( 18,2 \% ) obtiveram conceito 4 e somente um curso ficou com conceito 5, correspondendo a 0,9 \% dos cursos dessa região.
O Gráfico 5.1 apresenta a distribuição do Conceito Enade segundo Grande Região : poligonais mais à esquerda representam Grandes Regiões com pior distribuição, e poligonais mais à direita, Grandes Regiões com melhores conceitos.
A região Norte ( linha azul ) apresenta os cursos em IES com conceitos mais baixos ( localizada mais à esquerda no gráfico ), ao passo que a região Sudeste ( linha vermelha ) apresenta a distribuição com valores maiores ( localizada mais à direita no gráfico ).
Cursos em IES das regiões Centro-Oeste ( linha preta ), Nordeste ( linha verde ) e Sul ( linha roxa ) apresentam situações intermediárias. 5.2 CONCEITOS POR CATEGORIA ADMINISTRATIVA E POR GRANDE REGIÃO A Tabela 5.2 apresenta a distribuição dos cursos participantes do ENADE/2014 de Pedagogia ( Licenciatura ), por Categoria Administrativa, de acordo com os conceitos por eles alcançados, segundo as Grandes Regiões brasileiras.
Dos 1.126 cursos participantes, 337 ( 29,9 \% ) eram ministrados em Instituições Públicas e 789 ( 70,1 \% ), em Instituições Privadas.
De acordo com as informações da Tabela 5.2, em termos nacionais, entre Instituições Públicas e Privadas : das 35 IES que receberam o conceito 5, 12 eram Públicas e 23 Privadas.
Dos 337 cursos participantes de IES Públicas, o conceito 3 foi o valor modal, atribuído a 132 cursos ( 39,2 \% ).
Entre os demais cursos de Instituições Públicas participantes, 23 cursos ( 6,8 \% ) foram avaliados com conceito 1, 81 cursos ( 24,0 \% da categoria ) obtiveram conceito 2, 83 cursos ( 24,6 \% ) receberam conceito 4 e, como já comentado, 12 cursos ( 3,6 \% ) foram alocados ao conceito 5.
Nesta categoria, seis cursos ( 1,8 \% ) ficaram sem conceito ( SC ).
Na rede privada, o conceito modal foi 3, com 349 cursos dos 789 da categoria.
Entre os demais cursos participantes, 24 receberam conceito 1, e 195, conceito 2.
O conceito 4 foi atribuído a 181 cursos e o conceito 5, a 23 cursos.
Nesta Categoria Administrativa, outros 17 cursos ficaram sem conceito ( SC ).
Na análise por região, observa -se que, na região Norte, as Instituições Privadas participaram com 46 cursos ( 48,9 \% do total regional ), dos quais nenhum ficou com conceito 5 e dois ficaram sem conceito ( SC ).
O conceito modal para as Instituições Privadas na região foi 2, com 24 cursos, correspondendo a 52,2 \%.
E três Na região Norte, as Instituições Privadas participaram com 46 cursos ( 48,9 \% do total regional ), dos quais nenhum ficou com conceito 5 e dois ficaram sem conceito ( SC ).
O conceito modal para as Instituições Privadas na região foi 2, com 24 cursos, correspondendo a 52,2 \%.
E três cursos foram alocados no conceito 1, 15 no conceito 3 e outros dois no conceito 4.
As Instituições Públicas participaram com 48 cursos na região Norte ( 51,1 \% do total regional ), dos quais 22 obtiveram conceito 2, o valor modal, seguido por 16 cursos ( 33,3 \% ) que receberam conceito 3, cinco que receberam conceito 4 e quatro cursos que obteve conceito 1.
Apenas um curso ficou sem conceito ( SC ) e nenhum curso obteve conceito 5.
Na região Nordeste, a rede privada concentrou 100 dos 238 cursos participantes, o equivalente a 42,0 \% do total da região, a menor das proporções desta categoria na região.
Dentre os cursos oferecidos por IES Privadas no Nordeste, 40 cursos ficaram alocados ao conceito 3 ( conceito modal ), 32 cursos obtiveram conceito 2, 14 receberam conceito 4 e dez foram avaliados com conceito 1.
Três cursos ficaram sem conceito ( SC ) e apenas um curso foi avaliado com conceito 5.
As Instituições Públicas dessa região participaram com 138 cursos ( 58,0 \% ), a maior das proporções desta categoria na região, dos quais 61 ( 44,2 \% ) obtiveram conceito 3, o valor modal, seguido por 40 cursos que receberam conceito 2 ( 29,0 \% ).
Três cursos foram alocados ao conceito 5, 17 receberam conceito 1 e 16 foram avaliados com conceito 4.
Somente um curso ficou sem conceito ( SC ).
Na região Sudeste, a proporção de cursos da rede privada, 86,9 \%, foi mais elevada do que nas demais regiões brasileiras, correspondendo a 430 dos 495 cursos participantes.
Nesta categoria, na região Sudeste, o conceito modal foi 3 ( 209 cursos ).
Nesta combinação de Categoria Administrativa e Grande Região, nove cursos não receberam conceito ( SC ).
Os demais foram avaliados com conceito 1 ( cinco cursos ), conceito 2 ( 80 cursos ), conceito 4 ( 114 cursos ) e conceito 5 ( 13 cursos ).
Entre os 65 cursos oferecidos por Instituições Públicas na região Sudeste, o conceito modal foi 3 ( 23 cursos ), e os demais cursos receberam os conceitos 1 ( um curso ), 2 ( 11 cursos ), 4 ( 21 cursos ) e 5 ( oito cursos ).
Na região Sudeste, nesta categoria, somente um curso ficou sem conceito ( SC ).
As Instituições Privadas concentraram 142 dos 189 cursos participantes da região Sul, 75,1 \% do total regional.
Desses, 61 ficaram com conceito 3, o conceito modal.
Nesta combinação de Categoria Administrativa e Grande Região, 46 cursos ficaram com conceito 4, 24 cursos obtiveram conceito 2, enquanto oito cursos receberam conceito 5.
Nenhum curso foi associado ao conceito 1, enquanto três cursos ficaram sem conceito ( SC ).
As Instituições Públicas na região Sul participaram com 47 cursos ( 24,9 \% ) e o conceito modal foi 4, com 26 cursos.
Um curso foi avaliado com conceito 5, e 17 obtiveram conceito 3.
Nenhum curso foi avaliado com conceitos 1 ou 2.
Três cursos ficaram sem conceito ( SC ).
Na região Centro-Oeste, 71 dos 110 cursos participantes eram de Instituições Privadas ( 64,5 \% em termos regionais ).
Destes, 35 concentraram -se no conceito 2, conceito modal.
Dentre os demais, seis receberam conceito 1, 24 receberam conceito 3, cinco obtiveram conceito 4 e somente um curso foi avaliado com conceito 5.
Nenhum curso ficou sem conceito ( SC ).
Dos 39 cursos oferecidos por Instituições Públicas, 30 foram avaliados nos conceitos 3 e 4 ( 15 cursos cada, conceitos modais ), seguidos de oito cursos que obtiveram conceito 2, e um único curso ficou com conceito 1.
Nesta região, nenhum curso de IES pública ficou sem conceito ( SC ) ou recebeu conceito 5.
O Gráfico 5.2 apresenta a distribuição do Conceito Enade segundo a Categoria Administrativa da IES.
Os cursos oferecidos em IES Públicas ( linha azul ) apresentam uma poligonal praticamente igual àquela correspondente aos cursos oferecidos em IES Privadas ( linha verde ) e, consequentemente, uma distribuição de Conceitos Enade com valores muito próximos. 5.3 CONCEITOS POR ORGANIZAÇÃO ACADÊMICA E POR GRANDE REGIÃO Na Tabela 5.3, encontra -se a distribuição dos conceitos atribuídos aos cursos participantes do ENADE/2014 na Área de Pedagogia ( Licenciatura ), por Organização Acadêmica, segundo as Grandes Regiões brasileiras.
Dos 1.126 cursos de Pedagogia ( Licenciatura ) participantes, 441 eram oferecidos em Universidades, 101 cursos em Centros Universitários e 584 cursos oferecidos em Faculdades.
Esta distribuição corresponde a, respectivamente, 39,2 \%, 9,0 \% e 51,9 \% dos cursos.
De acordo com os dados apresentados, de todos os 35 cursos avaliados com conceito 5, 27 eram vinculados a Universidades.
Oito cursos oferecidos em Universidades ficaram sem conceito ( SC ).
Esse tipo de Organização Acadêmica teve o conceito 3 como modal, com 175 cursos ( 39,7 \% ).
Os demais cursos de Universidades avaliados receberam os conceitos 1 ( 23 cursos ), 2 ( 86 cursos ), 4 ( 122 cursos ) e conceito 5 ( 27 cursos, como já mencionado ).
Entre os cursos em Centros Universitários, o conceito modal foi 3, com 44 cursos ( 43,6 \% ).
Neste tipo de Organização Acadêmica, quatro cursos ficaram sem conceito ( SC ), e os outros cursos receberam os conceitos 2 ( 20 cursos ), 4 ( 32 cursos ) e 5 ( um curso ).
Nenhum curso recebeu conceito 1.
Nas Faculdades, 11 dos 584 cursos ficaram sem conceito ( SC ), e 262 ( 44,9 \% ) ficaram alocados no conceito 3, conceito modal, seguido de 170 cursos que obtiveram conceito 2 e de 110 cursos que receberam conceito 4.
Sete cursos foram avaliados com conceito 5 e 24 com conceito 1.
Considerando -se separadamente as regiões brasileiras, verifica -se que, na região Norte, as Universidades concentraram 47 dos 94 cursos participantes.
Desses, 21 cursos ficaram alocados ao conceito 2, conceito modal, seguidos de 16 cursos que obtiveram conceito 3, enquanto somente um curso não recebeu conceito ( SC ).
Cinco cursos foram avaliados com conceito 4 e quatro cursos obtiveram conceito 1.
Nenhum curso foi avaliado com conceito 5.
Os Centros Universitários da região Norte foram representados por somente quatro cursos, todos os quais receberam conceito 2.
As Faculdades participaram com 43 cursos na região Norte.
Desses, 21 foram avaliados com conceito 2, conceito modal, seguidos de 15 cursos que receberam conceito 3, de outros três cursos que obtiveram conceito 1 e de mais dois cursos que receberam conceito 4.
Apenas dois cursos ficaram sem conceito ( SC ) e nenhum curso obteve conceito 5.
Na região Nordeste, as Universidades participaram com 140 dos 238 cursos na Área de Pedagogia ( Licenciatura ).
O conceito modal foi 3, atribuído a 63 cursos, enquanto o conceito 4 foi atribuído a 16 cursos.
Três cursos obtiveram conceito 5, 40 cursos foram alocados no conceito 2 e 16 no conceito 1.
Dois cursos ficaram sem conceito ( SC ).
Os Centros Universitários contaram com somente quatro cursos participantes na região Nordeste, dois obtiveram conceito 3, conceito modal, e outros dois cursos, conceitos 2 e 4 ( um curso cada ).
As Faculdades foram representadas por 94 cursos na região Nordeste, dos quais 36 receberam conceito 3, conceito modal.
Os demais cursos foram avaliados com conceito 1 ( 11 cursos ), 2 ( 31 cursos ) e 5 ( 13 cursos ).
Dois cursos ficaram sem conceito ( SC ) e apenas um curso foi avaliado com conceito 5.
Na região Sudeste, as Universidades concentraram 127 dos 495 cursos da região.
Entre os cursos em Universidades, nessa região, o conceito modal foi 3 com 50 cursos, enquanto um curso ficou sem conceito ( SC ).
Os demais cursos receberam os conceitos 1 ( três cursos ), 2 ( 12 cursos ), 4 ( 42 cursos ) e 5 ( 19 cursos ).
Os Centros Universitários participaram com 66 cursos na região Sudeste, dos quais 31 obtiveram o conceito modal, 3, enquanto três cursos ficaram sem conceito ( SC ).
Nenhum curso recebeu conceitos 1 ou 5 nessa categoria e região.
Os demais receberam os conceitos 2 ( nove cursos ) e 4 ( 23 cursos ).
As Faculdades foram representadas por 302 cursos na região Sudeste dos quais 151 receberam o conceito 3, conceito modal.
Os demais se distribuíram nos conceitos 2 e 4 ( 70 cursos cada ), 1 ( três cursos ) e 5 ( dois cursos ).
Seis cursos ficaram sem conceito ( SC ) nessa categoria e região.
Dos 189 cursos da região Sul, 86 eram de Universidades, para os quais o conceito modal foi 4, atribuído a 44 cursos.
Nesse tipo de organização, nenhum dos cursos ficou com conceito 1.
Os demais receberam os conceitos 2 ( cinco cursos ), 3 ( 29 cursos ) e 5 ( quatro cursos ).
Quatro cursos ficaram sem conceito ( SC ).
Dos 16 cursos participantes de Centros Universitários da região Sul, sete obtiveram conceito 4, conceito modal.
Os demais cursos se distribuíram entre os conceitos 2 ( três cursos ), 3 ( quatro cursos ) e 5 ( um curso ).
Nenhum curso obteve conceito 1 e apenas um curso ficou sem conceito ( SC ).
Foram 87 cursos vinculados a Faculdades na região Sul, dos quais 45 receberam conceito 3, o conceito modal.
Quatro cursos receberam o conceito 5, enquanto 21 cursos ficaram com conceito 4 e 16 com conceito 2.
Somente um curso ficou sem conceito ( SC ) e nenhum curso foi avaliado com conceito 1 Na região Centro-Oeste, 41 dos 110 cursos eram de Universidades.
Nesse tipo de organização, nenhum curso ficou sem conceito ( SC ) ou com conceito 1.
O conceito modal foi 3, atribuído a 17 cursos, seguidos de 15 cursos que obtiveram conceito 4, de mais oito que receberam conceito 2 e de apenas um curso que foi alocado no conceito 5.
Os Centros Universitários da região Centro-Oeste contaram com 11 cursos, que obtiveram conceitos 3 ( sete cursos ), 2 ( três cursos ) e 4 ( um curso ).
Nenhum curso ficou sem conceito ( SC ) ou recebeu os conceitos restantes na categoria.
Dos 58 cursos oferecidos em Faculdades na região Centro-Oeste, 32 ficaram alocados no conceito 2, conceito modal.
Os demais cursos receberam os conceitos 1 ( sete cursos ), 3 ( 15 cursos ) e 4 ( quatro cursos ).
Nenhum curso ficou sem conceito ( SC ) ou com conceito 5.
O Gráfico 5.3 apresenta a distribuição do Conceito Enade segundo a Organização Acadêmica da IES.
Os cursos em Universidades ( linha azul ) e Centros Universitários ( linha verde ) apresentam uma poligonal mais à direita do que a de Faculdades, denotando uma melhor distribuição de Conceitos Enade.
No outro extremo, temos os cursos oferecidos em Faculdades ( linha vermelha ) com a poligonal mais à esquerda e, consequentemente, uma distribuição de Conceitos Enade com menores valores.

% matched lemmas: 
\end{textsample}
